\section*{Supplementary Material: Statistical Significance and Bootstrap Analysis}

To validate the robustness of our feature split ablation study, we conducted two rigorous statistical analyses: bootstrapping for confidence interval estimation and Edwards' continuity-corrected McNemar's test for predictive model comparison.

\subsection*{Bootstrap Confidence Intervals}
For each feature split (Terrain Only, Climate Only, Combined), we performed $B = 1,000$ bootstrap iterations. In each iteration, we sampled with replacement from the 2022 extreme flood test period dataset (representing the historic monsoon flood anomaly). For each bootstrap sample, we calculated the classification metrics: $F_1$-score, Inundation Intersection over Union (IoU), Precision, Recall, and Area Under the ROC Curve (AUC).

The 95\% confidence intervals were derived using the percentile method:
\begin{equation}
    CI_{95\%} = \left[ \theta^*_{0.025}, \theta^*_{0.975} \right]
\end{equation}
where $\theta^*_p$ represents the $p$-th percentile of the bootstrap distributions. The resulting intervals are compiled in Table~\ref{tab:significance_tests}. The complete overlap of the 95\% confidence intervals for $F_1$-score and IoU between the Terrain-only and Combined models fails to show a significant difference in their predictive performance.

\subsection*{Edwards' Continuity-Corrected McNemar's Test}
To determine if the differences in prediction vectors are statistically significant, we applied McNemar's test on the binary classification errors. McNemar's test is a non-parametric test for paired nominal data, evaluating the marginal homogeneity of a $2 \times 2$ contingency table.

The contingency table of discordant predictions is defined as:
\begin{center}
\begin{tabular}{|c|c|c|}
\hline
 & \textbf{Model 2 Correct} & \textbf{Model 2 Incorrect} \\
\hline
\textbf{Model 1 Correct} & $A$ & $B$ \\
\hline
\textbf{Model 1 Incorrect} & $C$ & $D$ \\
\hline
\end{tabular}
\end{center}
where Model 1 is the Terrain-only model and Model 2 is the Combined model.
\begin{itemize}
    \item $B$: Number of samples where Terrain-only is correct and Combined is incorrect ($B = 154$).
    \item $C$: Number of samples where Combined is correct and Terrain-only is incorrect ($C = 175$).
\end{itemize}

Edwards' continuity correction is applied to account for the discrete nature of the chi-square statistic:
\begin{equation}
    \chi^2 = \frac{(|B - C| - 0.5)^2}{B + C}
\end{equation}
Evaluating the discordant pairs ($B = 154$, $C = 175$):
\begin{equation}
    \chi^2 = \frac{(|154 - 175| - 0.5)^2}{154 + 175} = \frac{(20.5)^2}{329} = \frac{420.25}{329} \approx 1.2774
\end{equation}

For a chi-square distribution with 1 degree of freedom, the critical value for significance at $\alpha = 0.05$ is $\chi^2_{\text{crit}} = 3.841$. Since our calculated $\chi^2 = 1.2774 < 3.841$, we fail to reject the null hypothesis of marginal homogeneity. The corresponding p-value is:
\begin{equation}
    p = 0.2584 > 0.05
\end{equation}
This indicates that we cannot reject the null hypothesis that the predictions of the Terrain-only and Combined models differ in their error patterns, suggesting that the observed difference is consistent with sampling variability under limited statistical power rather than positive evidence of equivalence.

To quantify the magnitude of these differences, we calculate two standard effect size metrics for McNemar's test: the Odds Ratio ($OR$) and Cohen's $g$. The Odds Ratio is defined as:
\begin{equation}
    OR = \frac{B}{C} = \frac{154}{175} \approx 0.880
\end{equation}
where $OR$ represents the relative odds of errors between the two models (or inversely $C/B \approx 1.136$, representing that the Combined model has 1.136 times the odds of corrected predictions relative to Terrain-only). Cohen's $g$ measures the deviation of the discordant proportion from the null hypothesis of $0.5$:
\begin{equation}
    g = \left| \frac{B}{B + C} - 0.5 \right| = \left| \frac{154}{329} - 0.5 \right| = |0.4681 - 0.5| = 0.0319
\end{equation}
A Cohen's $g < 0.05$ is considered a negligible effect size, suggesting that the predictive difference between the Terrain-only and Combined models is both statistically non-significant and practically negligible.

For completeness, we also compared the Climate-only and Combined models ($B=356$ and $C=2357$ discordant pairs):
\begin{equation}
    OR = \frac{356}{2357} \approx 0.151, \quad g = \left| \frac{356}{2713} - 0.5 \right| \approx 0.3688
\end{equation}
yielding a highly significant difference ($\chi^2 = 1475.12$, $p < 0.0001$) and a very large effect size. Similarly, comparing the Climate-only and Terrain-only models ($B=479$ and $C=2459$ discordant pairs) yields:
\begin{equation}
    OR = \frac{479}{2459} \approx 0.195, \quad g = \left| \frac{479}{2938} - 0.5 \right| \approx 0.3370
\end{equation}
representing a highly significant difference ($\chi^2 = 1333.70$, $p < 0.0001$) and a large effect size, further supporting the dominant spatial control of geomorphology.


\subsection*{Spatial Autocorrelation and Point Independence Check}
To verify that our temporal training/validation point split does not suffer from spatial autocorrelation leakage, we conducted a spatial autocorrelation analysis using Global Moran's I on the 1,000 stratified sampling points. 

The spatial weight matrix was defined using an inverse-distance weighting scheme, and row-normalized. Statistical significance was evaluated using a permutation test with 999 random permutations of the target variables. 

The results (compiled in Table~\ref{tab:spatial_autocorrelation}) indicate that JRC water occurrence ($I = 0.209442$, $p = 0.0010$ for all points; $I = 0.209384$, $p = 0.0010$ for the validation subset) and land surface elevation ($I = 0.180902$, $p = 0.0010$) exhibit significant positive spatial autocorrelation, indicating a strong physical spatial structure in the delta. This spatial dependence poses a significant risk of spatial leakage if models exploit coordinate proximity during training. 

To verify that our disjoint point split and static terrain features successfully eliminate this leakage, we trained a Random Forest regressor on 800 training coordinates to predict occurrence based on elevation, slope, river distance, and coastal distance. We then computed Moran's I on the spatial residuals of the model predictions on the 200 unseen validation coordinates. Moran's I on validation-set residuals yields $I = 0.002022$ ($p = 0.8930$). To evaluate the effect size and significance of these Moran's I values, we compute Z-scores based on the permutation test distributions. For the validation-set residuals, the permutation test (Permutation Mean $E(I) = -0.005025$, standard deviation $\text{StdDev}(I) = 0.048321$) yields:
\begin{equation}
    Z = \frac{I - E(I)}{\text{StdDev}(I)} = \frac{0.002022 - (-0.005025)}{0.048321} \approx 0.146
\end{equation}
which is close to zero and suggests that no statistically significant spatial structure remains in the validation residuals. Conversely, for the baseline water occurrence across all points ($I = 0.209442$, Permutation Mean $E(I) = -0.001001$, standard deviation $\text{StdDev}(I) = 0.020015$), the Z-score is:
\begin{equation}
    Z = \frac{I - E(I)}{\text{StdDev}(I)} = \frac{0.209442 - (-0.001001)}{0.020015} \approx 10.51
\end{equation}
which represents an extremely large, highly significant spatial autocorrelation effect size before model fitting. This indicates that the model's static features successfully capture the spatial structure of surface water occurrence, leaving only spatially random errors on unseen coordinate geometries ($Z = 0.146$, $p = 0.8930$).

\begin{table}[t]
\caption{Global Moran's I spatial autocorrelation statistics and permutation-based significance testing results (999 permutations) for the stratified sampling points and Random Forest model spatial residuals on validation coordinates in the Indus Delta.}
\label{tab:spatial_autocorrelation}
\centering
\small
\begin{tabular}{p{4.2cm} c c p{6.5cm}}
\hline
\textbf{Variable / Residual} & \textbf{Moran's I} & \textbf{p-value} & \textbf{Spatial Interpretation} \\
\hline
Water Occurrence (All Points) & $0.209442$ & $0.0010$ & Significant positive spatial autocorrelation \\
Water Occurrence (Validation) & $0.209384$ & $0.0010$ & Significant positive spatial autocorrelation \\
Elevation (All Points) & $0.180902$ & $0.0010$ & Significant positive spatial autocorrelation \\
Random Forest Residuals (Validation) & $0.002022$ & $0.8930$ & No statistically significant spatial autocorrelation (spatially random) \\
\hline
\end{tabular}
\end{table}


\newpage
\section*{Supplementary Note S1: Physics-Constrained Hydrological State Space Model (PHSSM) Equations}

In Section 7.2 of the main text, we conceptually introduce a proposed Physics-Constrained Hydrological State Space Model (PHSSM) as a future modeling direction. This architecture couples a 1D temporal mass-balance state transition with a geomorphological spatial decoder. Here, we document the full parameterization and spatial decoder equations of the proposed architecture.

The temporal state transition tracks a latent equivalent water storage state $S_t$ (constrained by the mass-balance Equation 4 of the main text):
\begin{equation}
    S_t = S_{t-1} + P_t - ET_t - Q_t - G_t
\end{equation}
where $P_t$ is the CHIRPS monthly precipitation, and the dynamic fluxes---evapotranspiration $ET_t$, surface runoff $Q_t$, and groundwater infiltration $G_t$---are parameterized using differentiable Multi-Layer Perceptrons (MLPs):
\begin{align}
    ET_t &= \text{MLP}_{\text{ET}}(S_{t-1}, T_t, PET_t; \theta_{\text{ET}}) \\
    Q_t  &= \text{MLP}_{\text{Q}}(S_{t-1}, P_t; \theta_{\text{Q}}) \\
    G_t  &= \text{MLP}_{\text{G}}(S_{t-1}; \theta_{\text{G}})
\end{align}
where $T_t$ and $PET_t$ represent monthly temperature and potential evapotranspiration, and $\theta_{\text{ET}}$, $\theta_{\text{Q}}$, and $\theta_{\text{G}}$ represent the learnable parameters of the respective neural networks.

To resolve the sub-grid spatial representation collapse, the 1D temporal latent state $S_t$ is subsequently mapped to point-level water occurrence probabilities $p(y_{i,t} = 1)$ using a geomorphological spatial decoder. This decoder physically constrains the spatial inundation boundary using the static terrain features at each coordinate location $i$:
\begin{equation}
    p(y_{i,t} = 1) = \sigma \left( w_s S_t - w_z z_i - w_r D_{\text{river}, i} - w_c D_{\text{coast}, i} + b \right)
\end{equation}
where $z_i$ is the elevation, $D_{\text{river}, i}$ is the river distance, $D_{\text{coast}, i}$ is the coastal distance, $\sigma$ is the sigmoid activation function, $b$ is a learnable bias, and $w_s, w_z, w_r, w_c > 0$ are learnable positive weights (enforcing that higher storage increases inundation probability, while higher elevation and distance-decay functions decrease it).
